\documentclass[10pt,a4paper]{amsart}
\usepackage{amsmath,amssymb,amsthm,graphicx}
\usepackage{verbatim}
\newenvironment{reference}{\comment}{\endcomment}
%\newenvironment{reference}{}{}
\newenvironment{slogan}{\comment}{\endcomment}
\newenvironment{history}{\comment}{\endcomment}

% For commutative diagrams we use Xy-pic
\usepackage[all]{xy}

% We use 2cell for 2-commutative diagrams.
\xyoption{2cell}
% Theorem environments.
%
\theoremstyle{plain}
\newtheorem{theorem}[subsection]{Theorem}
\newtheorem{proposition}[subsection]{Proposition}
\newtheorem{lemma}[subsection]{Lemma}

\theoremstyle{definition}
\newtheorem{definition}[subsection]{Definition}
\newtheorem{example}[subsection]{Example}
\newtheorem{exercise}[subsection]{Exercise}
\newtheorem{situation}[subsection]{Situation}

\theoremstyle{remark}
\newtheorem{remark}[subsection]{Remark}
\newtheorem{remarks}[subsection]{Remarks}

\numberwithin{equation}{subsection}

% Macros
%
\def\lim{\mathop{\mathrm{lim}}\nolimits}
\def\colim{\mathop{\mathrm{colim}}\nolimits}
\def\Spec{\mathop{\mathrm{Spec}}}
\def\Hom{\mathop{\mathrm{Hom}}\nolimits}
\def\Ext{\mathop{\mathrm{Ext}}\nolimits}
\def\SheafHom{\mathop{\mathcal{H}\!\mathit{om}}\nolimits}
\def\SheafExt{\mathop{\mathcal{E}\!\mathit{xt}}\nolimits}
\def\Sch{\mathit{Sch}}
\def\Mor{\mathop{\mathrm{Mor}}\nolimits}
\def\Ob{\mathop{\mathrm{Ob}}\nolimits}
\def\Sh{\mathop{\mathit{Sh}}\nolimits}
\def\NL{\mathop{N\!L}\nolimits}
\def\CH{\mathop{\mathrm{CH}}\nolimits}
\def\proetale{{pro\text{-}\acute{e}tale}}
\def\etale{{\acute{e}tale}}
\def\QCoh{\mathit{QCoh}}
\def\Ker{\mathop{\mathrm{Ker}}}
\def\Im{\mathop{\mathrm{Im}}}
\def\Coker{\mathop{\mathrm{Coker}}}
\def\Coim{\mathop{\mathrm{Coim}}}

% Boxtimes
%
\DeclareMathSymbol{\boxtimes}{\mathbin}{AMSa}{"02}

%
% Macros for moduli stacks/spaces
%
\def\QCohstack{\mathcal{QC}\!\mathit{oh}}
\def\Cohstack{\mathcal{C}\!\mathit{oh}}
\def\Spacesstack{\mathcal{S}\!\mathit{paces}}
\def\Quotfunctor{\mathrm{Quot}}
\def\Hilbfunctor{\mathrm{Hilb}}
\def\Curvesstack{\mathcal{C}\!\mathit{urves}}
\def\Polarizedstack{\mathcal{P}\!\mathit{olarized}}
\def\Complexesstack{\mathcal{C}\!\mathit{omplexes}}
% \Pic is the operator that assigns to X its picard group, usage \Pic(X)
% \Picardstack_{X/B} denotes the Picard stack of X over B
% \Picardfunctor_{X/B} denotes the Picard functor of X over B
\def\Pic{\mathop{\mathrm{Pic}}\nolimits}
\def\Picardstack{\mathcal{P}\!\mathit{ic}}
\def\Picardfunctor{\mathrm{Pic}}
\def\Deformationcategory{\mathcal{D}\!\mathit{ef}}
\usepackage{cleveref}
\crefname{originallemma}{original lemma}{original lemmas}
\Crefname{originallemma}{Original lemma}{Original lemmas}
\makeatletter
\newlabel{lemma-flat-change-of-rings}{{05V4}{source}{}{sourcepointer-1}{}}
\newlabel{lemma-flat-change-of-rings@cref}{{[originallemma][0][]05V4}{[1][0][]source}{}{}{}}
\makeatother
\crefname{originallemma}{original lemma}{original lemmas}
\Crefname{originallemma}{Original lemma}{Original lemmas}
\makeatletter
\newlabel{lemma-flat-tor-zero}{{03EX}{source}{}{sourcepointer-2}{}}
\newlabel{lemma-flat-tor-zero@cref}{{[originallemma][0][]03EX}{[1][0][]source}{}{}{}}
\makeatother
\title{Flatness of a module over a square-zero thickening is equivalent to flatness of its reduction plus injectivity of the conormal multiplication}
\author{The Stacks Project Authors}
\date{stacks-project:sites-modules.tex:local-snapshot-20261012; GNU FDL 1.2 or later}
\begin{document}
\maketitle
\noindent\textit{Editorial note.} All delivered mathematical text is the author's own, copied verbatim from the pinned Stacks Project chapter file sites-modules.tex; no statement or proof step was rewritten, reordered or removed.\par\medskip
\noindent\textit{Editorial note.} The proof contains a commutative diagram typeset with Xy-pic; the retained two-slice preamble of this source supplies the native Xy-pic setup, so it compiles as in the source.\par\medskip
\noindent\textit{Editorial note.} The counted span cites the two cited lemmas and one section, all recorded in assembly.external\_references with canonical tags as navigation pointers.\par\medskip
\noindent\textit{Allowed standard prerequisites (editorial).}\begin{itemize}
\item Flatness is stable under change of rings and detectable by vanishing of Tor. Stacks Project, Modules on Sites, Lemmas lemma-flat-change-of-rings (tag 09B7) and lemma-flat-tor-zero (tag 09B8). Supplies the forward direction of the equivalence.
\item Tensor products of modules annihilated by a square-zero ideal. Stacks Project, Modules on Sites, section on modules over ringed topoi. Justifies the identification K tensor\_\{O'\} F' = K tensor\_O F used throughout the converse.
\end{itemize}\medskip
\ifdefined\hypertarget\hypertarget{sourcepointer-1}{}\fi
\noindent\textit{Source reference.} Original lemma 05V4: sources/stacks-chow/tags.txt, source lines 7538--7538. Cross-chapter standard prerequisite cited by the author inside the counted span; it is used as a pointer and its own proof is not part of this item.\par\medskip
\ifdefined\hypertarget\hypertarget{sourcepointer-2}{}\fi
\noindent\textit{Source reference.} Original lemma 03EX: sources/stacks-chow/tags.txt, source lines 4472--4472. Cross-chapter standard prerequisite cited by the author inside the counted span; it is used as a pointer and its own proof is not part of this item.\par\medskip
\begin{lemma}
\label{lemma-flat-over-thickening}
Let $\mathcal{C}$ be a site. Let $\mathcal{O}' \to \mathcal{O}$
be a surjection of sheaves of rings whose kernel $\mathcal{I}$ is
an ideal of square zero. Let $\mathcal{F}'$ be an $\mathcal{O}'$-module
and set $\mathcal{F} = \mathcal{F}'/\mathcal{I}\mathcal{F}'$.
The following are equivalent
\begin{enumerate}
\item $\mathcal{F}'$ is a flat $\mathcal{O}'$-module, and
\item $\mathcal{F}$ is a flat $\mathcal{O}$-module and
$\mathcal{I} \otimes_\mathcal{O} \mathcal{F} \to \mathcal{F}'$
is injective.
\end{enumerate}
\end{lemma}

\begin{proof}
If (1) holds, then
$\mathcal{F} = \mathcal{F}' \otimes_{\mathcal{O}'} \mathcal{O}$ is
flat over $\mathcal{O}$ by
Lemma \ref{lemma-flat-change-of-rings}
and we see the map
$\mathcal{I} \otimes_\mathcal{O} \mathcal{F} \to \mathcal{F}'$
is injective by applying $- \otimes_{\mathcal{O}'} \mathcal{F}'$
to the exact sequence
$0 \to \mathcal{I} \to \mathcal{O}' \to \mathcal{O} \to 0$, see
Lemma \ref{lemma-flat-tor-zero}.
Assume (2). In the rest of the proof we will use without further mention
that $\mathcal{K} \otimes_{\mathcal{O}'} \mathcal{F}' =
\mathcal{K} \otimes_\mathcal{O} \mathcal{F}$
for any $\mathcal{O}'$-module $\mathcal{K}$ annihilated by $\mathcal{I}$.
Let $\alpha : \mathcal{G}' \to \mathcal{H}'$
be an injective map of $\mathcal{O}'$-modules. Let
$\mathcal{G} \subset \mathcal{G}'$, resp.\ $\mathcal{H} \subset \mathcal{H}'$
be the subsheaf of sections annihilated by $\mathcal{I}$.
Consider the diagram
$$
\xymatrix{
\mathcal{G} \otimes_{\mathcal{O}'} \mathcal{F}' \ar[r] \ar[d] &
\mathcal{G}' \otimes_{\mathcal{O}'} \mathcal{F}' \ar[r] \ar[d] &
\mathcal{G}'/\mathcal{G} \otimes_{\mathcal{O}'} \mathcal{F}' \ar[r] \ar[d] &
0 \\
\mathcal{H} \otimes_{\mathcal{O}'} \mathcal{F}' \ar[r] &
\mathcal{H}' \otimes_{\mathcal{O}'} \mathcal{F}' \ar[r] &
\mathcal{H}'/\mathcal{H} \otimes_{\mathcal{O}'} \mathcal{F}' \ar[r] & 0
}
$$
Note that $\mathcal{G}'/\mathcal{G}$ and $\mathcal{H}'/\mathcal{H}$
are annihilated by $\mathcal{I}$ and that
$\mathcal{G}'/\mathcal{G} \to \mathcal{H}'/\mathcal{H}$ is injective.
Thus the right vertical arrow is injective as $\mathcal{F}$ is flat
over $\mathcal{O}$. The same is true for the left vertical arrow.
Hence the middle vertical arrow is injective and $\mathcal{F}'$ is flat.
\end{proof}

\end{document}
